\subsectionlink{de Sitter極限とゴースト凝縮}

前節では，$c_s \to 0$ の極限では理論がますます強結合になるため，この極限が病的であることを見た．しかし，これまでの議論では，ユニタリーゲージのラグランジアン（\ref{eq:actiontad}）に含まれる高階微分演算子
\be
\int d^4 x \sqrt{-g} \left(-\frac{\bar M_2(t)^2}{2} \delta K^\mu {}_\mu {}^2
-\frac{\bar M_3(t)^2}{2} \delta K^\mu {}_\nu \delta K^\nu {}_\mu
\right) \;.
\ee
を無視していた．これらの演算子は，Goldstone作用において
\begin{equation}
\label{di_pi2}  \int  \! d^4x \: \sqrt{-g} \left[ - \frac{\bar M^2}{2} \, \frac{1}{a^4}(\di_i^2 \pi)^2
   \right] \; ,
\end{equation}
という形の空間的運動項を生む．ここで，$\bar M^2 = \bar M_2^2+  \bar M_3^2$ である．
この空間的運動項によって，$c_s \to 0$ の極限でもGoldstoneモードは伝播する．したがって，$\dot H = 0$ の極限で，一般的なラグランジアンを考察することは興味深い．この極限では重力背景は厳密なde Sitter時空となり，$c_s=0$ が成り立つ．この場合，$H$ は時間に依存しないため，理論に追加の対称性を課すことができる．すなわち，ユニタリーゲージのラグランジアンに含まれるすべての係数を時間に依存しないものとする．Goldstoneモード $\pi$ を再導入する手順（\ref{Smixed}）を振り返ると，この対称性は微分を伴わない $\pi$ への依存をすべて禁止することがわかる．したがって，Goldstone作用は $\pi$ のシフト
\be
\pi(\vec x,t) \to \pi(\vec x,t) + {\rm const.}
\ee
のもとで不変である．これがゴースト凝縮 \cite{Arkani-Hamed:2003pdi} の極限であり，Goldstoneモードは非相対論的な分散関係 $\omega \propto k^2$ を持つ．より一般には，空間的運動項 $c_s^2
(\nabla\pi)^2$ と高階微分項 $(\nabla^2\pi)^2$ の両方が存在する中間的な状況を考えることができる．理論の予言は，凍結が起こるエネルギー $\omega
\sim H$ においてどちらの項が支配的であるかによって，大きく変化する \cite{Arkani-Hamed:2003juy,Senatore:2004rj}．

これまでの模型と同様に，Goldstoneモードと重力の混合を無視できるエネルギー領域を求めなければならない．この領域を見積もる簡単な方法の一つは，$\delta K$ 演算子が含む次のような項を見ることである．
\be
\delta K_{ij} \supset (\partial_i\partial_j \pi +\partial_i g_{0j}) \;.
\ee
正準規格化を行うと，$k \gtrsim M_2^2/M_{\rm Pl}$ では重力との混合を無視できることがわかる．Goldstoneモードの分散関係は $\omega^2 = \bar M^2/M_2^4 \cdot k^4$ の形なので，混合が重要になるエネルギーの上限 $E_{\rm mix}$ は
$E_{\rm mix} \simeq \bar M M_2^2/M_{\rm Pl}^2$ であることがわかる \cite{Arkani-Hamed:2003pdi}．
このスケールは背景の曲率とは無関係であることに注意しよう．これは，この例の非常に顕著な特徴である．通常，重力との混合は背景のエネルギー・運動量テンソル，したがって時空の曲率に関係しており，系が空間を強く湾曲させるほど，重力との混合も強くなる．一方，この場合には，平坦なMinkowski背景上でも混合が重要になる．これが本来の意味での重力の修正と呼ばれるものである．すなわち，たとえば源が生成するNewtonポテンシャルのような重力が，曲率スケールよりもはるかに小さなスケールで修正される．この修正重力模型とその豊かな現象論は，文献 \cite{Arkani-Hamed:2003pdi} で研究されている\footnote{音速が低下した模型の場合にも，重力との混合のスケールは地平線よりパラメトリックにはるかに小さくなりうる．$\epsilon/c_s^2 \ll 1$ であれば十分である．この場合，その模型は重力を修正する一つの方法と見なすことができる．ただし，空間的な高階微分項を考慮せずに $\dot H = 0$ の極限をとることはできない．そうしなければ，スカラー・モードは伝播しなくなる．}．

ここではインフレーションに関心があるので，逆の極限 $H \gg E_{\rm mix}$ に注目する．この極限では混合を無視でき，$\pi$ のラグランジアンに集中できる．ゴースト・インフレーションの主な特徴を簡単に説明しよう．詳細については文献
\cite{Arkani-Hamed:2003pdi,Arkani-Hamed:2003juy,Senatore:2004rj} を参照されたい．これらの文献では，本論文で提示したものに非常に近いアプローチで理論が研究されている．興味深い特徴の大半は，各演算子のエネルギーに関するスケーリングを見ることで理解できる．非相対論的な分散関係 $\omega \propto k^2$ のため，演算子のエネルギーに関するスケーリングは，Lorentz不変な場合のようにその質量次元とは一致しない．
エネルギーを因子 $s$ で再スケールする変換 $E\rightarrow s E$（時間の再スケール $t\rightarrow s^{-1} t$ と等価である）には，運動量の変換 $k\rightarrow s^{1/2}k$（空間座標では $x\rightarrow
s^{-1/2}x$）が伴わなければならない．Goldstoneモードの二次の作用は
\be
\label{eq:kinghost}
\int d^4 x \left[ 2 M_2^4 \dot\pi^2 - \frac{\bar M^2}{2} \,
  \frac{1}{a^4}(\di_i^2 \pi)^2 \right]
\ee
という形なので，二次の作用を不変に保つには，$\pi$ にスケーリング次元 $1/4$ を割り当てなければならない．すなわち，
\begin{equation}
\pi\rightarrow s^{1/4}\pi \, .
\end{equation}
この規則を用いると，運動項（\ref{eq:kinghost}）以外の，許されるすべてのGoldstone演算子が非関連であることを容易に確かめられる．{\em すなわち}，これらは正のスケーリング次元を持ち，エネルギーを下げるほど重要性を失っていく．このことは，この理論が有効場理論として意味を持つことを示している．特に，高階時間微分を含む運動演算子 $\ddot \pi^2$ は，一見すると $(\nabla^2\pi)^2$ と同程度に重要に見え，理論にゴーストが存在することを記述するが，次元が 2 であり，低エネルギーでは無視できる．問題に単一のスケール $M$ があり，$M \simeq M_2 \simeq \bar M$ であると仮定すれば，このスケールが有効場理論による記述のエネルギー・カットオフも定める．

$\pi$ のスケーリング次元からは，ゴースト・インフレーションで生成される摂動のスペクトルも見積もることができる．$\pi$ の次元は，量子揺らぎの振幅がエネルギーとともにどのように変化するかを教えてくれる．カットオフのスケール $\Lambda \simeq M$ では，正準規格化されたGoldstone場 $\pi_c \simeq M^2 \pi $ の量子揺らぎは，カットオフと同程度の大きさ $\delta \pi_c (M)\sim M$ である．エネルギーを下げることで，凍結時 $E \sim H$ の量子揺らぎを見積もることができる．標準的な場合，スカラー場の次元は 1 であり，凍結時の揺らぎは $H$ 程度となるが，ここでは
\begin{equation}
\delta\pi_c(H) \sim \delta\pi_c(M)
\left(\frac{H}{M}\right)^{1/4}\sim \left(H M^3\right)^{1/4} \;.
\end{equation}
となる．スケール $H$ での量子揺らぎは，Lorentz不変な分散関係を持つスカラー場と比べて大幅に増大している．したがって，$\zeta$ のスペクトルは \cite{Arkani-Hamed:2003juy}
\begingroup\addtolength{\belowdisplayskip}{-2pt}
\be
\label{eq:spectrumzetaghost}
\langle \zeta(\vec k_1) \zeta(\vec k_2)\rangle \sim (2 \pi)^3 \delta(\vec
k_1 + \vec k_2)  \frac{H^2}{M^4} (H M^3)^{1/2} \frac{1}{k_1^3} \;.
\ee
\endgroup
で与えられる．スペクトルが正しい規格化を持つためには，$(H/M)^{5/4} \simeq
10^{-5}$ が必要である．

非Gaussian性は，次元が最も低い演算子によって支配される．演算子 $M_2^4
\dot\pi (\nabla\pi)^2$ は，ユニタリーゲージの演算子 $M_2^4
(g^{00}+1)^2$ から生じ，次元 $1/4$ を持つ最も非関連性の弱い演算子であることが容易にわかる．
カットオフのスケール $M$ では，理論は強結合である．三次の演算子は次元 $1/4$ を持つため，$H$ 程度のエネルギーでは，$(H/M)^{1/4}$ 程度の非Gaussian性を生み，これはパラメトリックには $P_\zeta^{1/5}$ 程度である．同じ結果は，前節で用いたアプローチ，{\em すなわち}凍結時に相互作用項と自由作用を比較することによっても得られる．
\begin{equation}
\label{eq:NGghost}
\frac{{\cal L}_{\dot\pi(\nabla \pi)^2}}{{\cal L}_{2}}\sim
\frac{M^4 H \pi (H M) \pi^2}{M^4 H^2 \pi^2}=M \pi=\frac{M}{H}
\zeta\sim\left(\frac{H}{M}\right)^{1/4} \, .
\end{equation}
最も関連性の高い演算子の次元が非常に低いため，非Gaussian性の大きさは標準的なスローロールの場合と比べて極めて大きい．具体的な計算
\cite{Arkani-Hamed:2003juy} では，効果は素朴な見積もりよりいくぶん小さくなり，既存の実験的な制限 \cite{Creminelli:2006rz} と同程度になる．

これまでの議論では，ユニタリーゲージの演算子
\begin{equation}
\label{eq:leooperator}
\int \! d^4 x \: \sqrt{- g} \:
\left(-\frac{\bar M_1(t)^3}{2} (g^{00}+1)\delta K^\mu {}_\mu \right)\;.
\end{equation}
を無視してきた．この演算子は時間反転のもとで奇であるため，その係数をゼロにすることは無矛盾である．この項が存在すれば，Goldstoneラグランジアンには，次元 $1/4$ を持つ二つ目の演算子が，$\nabla^2\pi (\nabla\pi)^2$ という形で現れる．その 3 点関数への寄与は，$\dot\pi(\nabla\pi)^2$ の寄与と同程度となる．ユニタリーゲージの演算子（\ref{eq:leooperator}）は，これから論じるように，$\pi$ の二次のラグランジアンにも寄与する．

\subsubsectionlink{ゴースト凝縮を伴わないde Sitter極限}
本節では，演算子（\ref{eq:leooperator}）が $\pi$ の二次の作用に及ぼす効果を調べる．この演算子の係数が十分に大きければ，新しいde Sitter極限が得られることを見る．この極限では，凍結時の分散関係が $\omega^2
\propto k^2$ の形となり，ゴースト凝縮の振る舞い $\omega^2
\propto k^4$ に取って代わる．

簡単のため，$\bar M_1$ は時間に依存しないとする．Goldstoneモードを再導入すると，$-\bar M_1^3 \dot\pi\nabla^2\pi/a^2$ という形の三階微分項が得られる（\footnote{この演算子は，$\dot\pi^2$ に対して，$H$ に比例する寄与も与える．この寄与は $M_2^4 \dot\pi^2$ に比べて小さいと仮定する．Minkowski時空では，ここで考えている演算子を $\phi \to -\phi$ 対称性によって禁止できる．この対称性は，ユニタリーゲージでは時間反転と等価である \cite{Arkani-Hamed:2003pdi}．de Sitter背景では，この対称性は計量によって破られるため，この演算子の係数をゼロにすることはできない．}）．スケール因子 $a(t)$ と計量の行列式に時間依存性がなければ，これは時間に関する全微分となる．部分積分すると，標準的な二階微分の空間的運動項
\be
\label{eq:byparts}
-\int d^4 x \sqrt{-g} \,\frac{\bar M_1^3 H}{2}
\left(\frac{\partial_i}{a}\pi\right)^2 \;.
\ee
が得られる．厳密なde Sitter極限 $\dot H =0$ で，$M_2 \sim \bar
M_1 \sim M$ とすると，この演算子は $\omega^2 = c_s^2 k^2$ という形の分散関係を与え，音速は小さくなる\footnote{この模型では，重力との混合はこれまでの場合とかなり異なる．その理由は，式~（\ref{eq:byparts}）を得るために時間微分を部分積分しているため，Goldstoneの項が，重力との混合を記述する項と同じ数の微分を含むことにある．このことは，高エネルギーに向かっても混合の重要性が低下していかないことを意味する．一方，混合が常に無視できるように模型のパラメータを選ぶことは可能である．具体的な計算については文献 \cite{Creminelli:2006rz} を参照されたい．}．
\begin{equation}
c^2_s=\frac{H}{M}\ll 1 \;.
\end{equation}
これが成り立つのは，高階微分演算子 $\delta K^\mu {}_\mu {}^2$ と
$\delta K^\mu {}_\nu \delta K^\nu {}_\mu$ が劣勢である場合に限られる．これらも同じ質量スケールで特徴づけられ，$\bar M_2 \sim \bar
M_3 \sim M$ であると仮定すると，分散関係には二つの寄与が生じる．
\be
\omega^2 \sim \frac{H}{M} k^2 + \frac{k^4}{M^2} \;.
\ee
この二つの空間的運動項は，凍結時 $\omega \sim
H$ には同程度となる．一方，$k^4$ の寄与がいくらか抑制されていれば，凍結時には無視でき，したがってインフレーションの予言にとっても重要ではなくなる．この極限では，厳密なde Sitter背景を持つ一方で，凍結時には $\omega^2 \propto k^2$ の分散関係を持つ，新しい種類のゴースト・インフレーションが得られる．

他の模型に対して行った手順に従えば，スペクトルの規格化と，3 点関数の非Gaussian性の見積もりを容易に得ることができる．
\be
\label{eq:spectrumzetaleo}
\langle \zeta(\vec k_1) \zeta(\vec k_2)\rangle \sim (2 \pi)^3 \delta(\vec
k_1 + \vec k_2)  \frac{H^4}{M^4} \frac1{c_s^3} \frac{1}{k_1^3}  \sim (2 \pi)^3 \delta(\vec
k_1 + \vec k_2)  \left(\frac{H}{M}\right)^{5/2} \frac{1}{k_1^3} \;.
\ee
\be
\frac{{\cal L}_{\dot\pi(\nabla \pi)^2}}{{\cal L}_{2}}\sim
\frac1{c_s^2} H \pi \sim\left(\frac{H}{M}\right)^{1/4} \, .
\ee
Goldstone演算子 $(\nabla\pi)^2\nabla^2\pi$ からも，同程度の寄与が生じる．当然ながら，上の見積もりは，ゴースト凝縮の場合に得た式（\ref{eq:spectrumzetaghost}）および（\ref{eq:NGghost}）と同じである．実際，すでに論じたように，すべての演算子を同じスケールにとると，凍結時には $k^2$ と $k^4$ の空間的運動項から同程度の寄与が得られる．したがって，二つの寄与のうち一方だけが存在すると仮定した場合にも，同様の予言が得られると期待される．

凍結時に $(g^{00}+1)\delta K^\mu {}_\mu $ が支配する場合と，$\delta K^\mu {}_\mu {}^2$ および $\delta K^\mu {}_\nu \delta K^\nu
{}_\mu$ が支配する場合という，二つの異なるde Sitter極限が見つかった．そこで，他の可能性もあるのではないかと考えるかもしれない．何らかの理由で，これら二つの空間的運動項がともに抑制され，より高次で初めて主要な演算子が現れることを想像できる．この場合，最終的には
\begin{equation}
\omega^2\sim k^{2n} \quad n \geq 3 \;.
\end{equation}
という形の分散関係が得られる．しかし，このようなことは起こりえないと容易にわかる．なぜなら，その理論は有効場理論として意味を持たなくなるからである．ゴースト凝縮で用いたのと同じ論理に従うと，演算子 $\pi$ のスケーリング次元は
\begin{equation}
\pi\rightarrow s^{-\frac12+\frac{3}{2n}} \pi \, .
\end{equation}
となる．このことから，演算子 $\dot\pi(\nabla\pi)^2$ は，対称性によって時間的運動項 $\dot\pi^2$ と結びついており，次元
$(7-3n)/(2n)$ を持つ．$n \geq 3$ では，この演算子は低エネルギーで強結合となるため，有効場理論は意味を持たない．

\sectionlink{結論}
インフレーションを検証するための実験的な取り組みが進行し，さまざまな模型が数多く提案されていることを考えると，インフレーションの最も一般的な理論を特徴づけることは非常に重要である．本論文では，新しい観点を採用した．一般的なラグランジアンを書き下し，インフレーションを起こす解のまわりの揺らぎを調べる代わりに，準de Sitter背景のまわりの揺らぎの有効理論を直接記述する．この理論では，空間微分同相変換は顕在的であり，時間微分同相変換は非線形に実現される．微分の最低次では，最も一般的な作用を次の形で書けることを示した．
\be
S &=& \int  \! d^4 x \; \sqrt{- g} \Big[ \frac12 M_{\rm Pl}^2 R + M_{\rm Pl}^2 \dot H
g^{00} - M_{\rm Pl}^2 (3 H^2 + \dot H) +
 \nonumber \\ && \frac{1}{2!}M_2(t)^4(g^{00}+1)^2+\frac{1}{3!}M_3(t)^4 (g^{00}+1)^3+
 \nonumber \\
&& - \frac{\bar M_1(t)^3}{2} (g^{00}+1)\delta K^\mu {}_\mu
-\frac{\bar M_2(t)^2}{2} \delta K^\mu {}_\mu {}^2
-\frac{\bar M_3(t)^2}{2} \delta K^\mu {}_\nu \delta K^\nu {}_\mu + ...
\Big] \; .
\ee

宇宙論的な相関関数は，Hubbleパラメータ $H$ 程度のスケールでこの有効場理論を検証する．このアプローチでは，対称性の役割がはるかに明確になる．どの特徴がインフレーションを起こす背景解から導かれるかを具体的に見ることができる．特に，$\dot H<0$，$\dot H = 0$，$\dot H > 0$ の各場合で，振る舞いが大きく異なることがわかる．これは，演算子 $g^{00}$ の係数が $\dot H$ によって固定されるためである．この観点から，本アプローチは，インフレーション，修正重力理論，およびバウンス・模型 \cite{Creminelli:2006xe,Buchbinder:2007ad,Creminelli:2007aq} のようにヌル・エネルギー条件を破る理論（宇宙論の文脈では $\dot H >0$ と等価）相互の関係をより明確にする．対称性の役割を示すもう一つの例は，音速の低下と非Gaussian性の増大の結びつきである．両者は同じ演算子 $M_2(t)^4(g^{00}+1)^2$ から生じるため，時間微分同相変換の非線形実現によって関連づけられている．

最も標準的なスローロールのシナリオからのあらゆる逸脱は，素粒子物理の標準模型の研究におけるのと同様に，高次演算子の大きさに体系的に組み込まれている．さらに，すべての単一場模型が共通の枠組みに統一されるため，たとえば文献 \cite{Cheung:2007sv,Arkani-Hamed:2007ryv} で行われたように，具体的な実現の仕方に依存しない一般的な結論を導くことが可能になる．

本定式化の拡張として考えられるものは容易に思い浮かぶ．同じ方針に沿って，（単一場の）クインテッセンスの最も一般的な理論を調べ，その現象論的な帰結を求めることは興味深いだろう．有効理論を $H$ 程度のスケールで調べるインフレーションとは異なり，この場合に関心があるのは，地平線内での摂動の力学である．放射優勢期や物質優勢期のような流体における揺らぎの研究 \cite{Dubovsky:2005xd} に，本アプローチを用いることも興味深い．最後に，追加の場を導入して，多場インフレーション・模型を研究することも容易であるはずだ．

\phantomsection\section*{\hyperlink{toc}{謝辞}}\addcontentsline{toc}{section}{謝辞}
多くの有益な議論をしていただいた N.~Arkani-Hamed，L.~Boubekeur，S.~Dubovsky，
A.~Guth，F.~Vernizzi，J.~Wacker，M.~Zaldarriaga，そして特に A.~Nicolis に感謝する．
A. Liam Fitzpatrick は NSF，Jared Kaplan は Hertz と NSF のフェローシップの支援を受けている．

\appendix

\section*{\hyperlink{toc}{付録}}

\sectionlink{ユニタリーゲージにおける最も一般的なラグランジアン\label{app:generic}}
ユニタリーゲージにおける最も一般的なラグランジアンを書き下すための規則を調べよう．空間微分同相変換のもとでのみ不変な理論には，関数 $\tilde t(x)$（時間的な勾配を持つ）が定める特別な時空の葉層分解があり，この関数が時間微分同相変換を非線形に実現する．
たとえば，時間発展するスカラー場によって対称性が破られている場合，$\tilde t$ が一定の面は，スカラー場の値が一定の面でもある．ユニタリーゲージとは，時間座標 $t$ が $\tilde t$ と一致するように選んだゲージであり，このため追加の自由度 $\tilde t$ は作用に陽には現れない．したがって，次のようなさまざまな項を構成できる．
\begin{enumerate}
\item すべての微分同相変換のもとで不変な項．これらは，Riemannテンソル $R_{\mu\nu\rho\sigma}$ とその共変微分からなる多項式を，スカラーになるように縮約したものにほかならない\footnote{添字の縮約には，計量と完全反対称テンソル $(-g)^{-1/2} \epsilon^{\mu\nu\rho\sigma}$ を用いることができる．}．
\item $\tilde t$ の一般的な関数は，ユニタリーゲージでは $f(t)$ になる．したがって，ラグランジアンの任意の項の前に，時間の一般的な関数を自由に掛けることができる．
\item 勾配 $\partial_\mu \tilde t$ は，ユニタリーゲージでは $\delta_\mu^0$ になる．したがって，すべてのテンソルにおいて，上付きの $0$ 添字を縮約せずに残すことが常にできる．たとえば，ユニタリーゲージのラグランジアンでは，$g^{00}$（およびその関数）や，Ricciテンソルの成分 $R^{00}$ を用いることができる．
\item $\tilde t$ が一定の面に垂直な単位ベクトルを定義すると便利である．
\be
n_\mu = \frac{\partial_\mu\tilde t}{\sqrt{-g^{\mu\nu}\partial_\mu \tilde t \partial_\nu\tilde t}}  \;.
\ee
これにより，$\tilde t$ が一定の面上の誘導空間計量 $h_{\mu\nu} \equiv g_{\mu\nu} + n_\mu n_\nu$ を定義できる．あらゆるテンソルは，$h_{\mu\nu}$ を用いてこれらの面上に射影できる．特に，誘導された三次元計量のRiemannテンソル $^{(3)} R_{\alpha\beta\gamma\delta}$ と，三次元計量に関する共変微分を作用に用いることができる．
\item $\partial_\mu\tilde t$ の共変微分から，さらに別の可能性が得られる．同じこととして，$n_\mu$ の共変微分を考えてもよいことに注意しよう．規格化因子に微分が作用しても，$\partial_\mu g^{00}$ のような項が得られるだけであり，これらはそれ自体で共変で，ユニタリーゲージのラグランジアンに用いることができる．$n_\mu$ の共変微分を $\tilde t$ が一定の面上に射影すると，これらの面の外的曲率が得られる．
\be
K_{\mu\nu} \equiv h_\mu^\sigma \nabla_\sigma n_\nu \;.
\ee
添字 $\nu$ はすでに面上に射影されている．これは，$n^\nu \nabla_\sigma n_\nu = \frac12 \nabla_\sigma (n^\nu n_\nu) =0$ であるためである．面に垂直な方向の $n_\nu $ の共変微分は，
\be
n^\sigma \nabla_\sigma n_\nu = - \frac12 (-g^{00})^{-1} h^\mu_\nu \partial_\mu(-g^{00})
\ee
と書き換えられるため，新しい項を生じない．したがって，$n_\mu$ のすべての共変微分は，外的曲率 $K_{\mu\nu}$（およびその共変微分）と $g^{00}$ の微分を用いて書くことができる．

\item 誘導された三次元計量のRiemannテンソルと外的曲率を同時に用いるのは冗長であることに注意しよう．$^{(3)} R_{\alpha\beta\gamma\delta}$ は，Gauss--Codazzi関係を用いて \cite{Wald:1984rg}
\be
^{(3)} R_{\alpha\beta\gamma\delta}= h^\mu_\alpha h^\nu_\beta h^\rho_\gamma h^\sigma_\delta R_{\mu\nu\rho\sigma} -K_{\alpha\gamma}K_{\beta\delta} +K_{\beta\gamma}K_{\alpha\delta} \;.
\ee
と書き換えられるからである．したがって，三次元Riemannテンソルを用いる必要はまったくない．誘導計量 $h_{\alpha\beta}$ を陽に用いることも避けられる．四次元計量と $n_\mu$ で書き表せば，すでに上で論じた項だけが得られる\footnote{誘導計量の行列式は，四次元の全計量の行列式と $h = g^{00} \cdot g$ で関係づけられること，また，三次元の完全反対称テンソルは，四次元のものを用いて $h^{-1/2} \epsilon^{ijk} = (-g)^{-1/2} (-g^{00})^{-1/2} \epsilon^{0ijk}$ と書き換えられることに注意しよう．}．最後に，誘導された三次元計量に関する共変微分を用いることも避けられる．射影されたテンソルの三次元共変微分は，四次元共変微分を射影することで得られる \cite{Wald:1984rg}．
\end{enumerate}

以上より，ユニタリーゲージにおける最も一般的な作用は
\be
S = \int d^4 x \sqrt{-g} \; F(R_{\mu\nu\rho\sigma}, g^{00}, K_{\mu\nu},\nabla_\mu,t) \;,
\ee
で与えられると結論できる．ここで，関数 $F$ の内部で縮約されずに残る添字は，すべて上付きの 0 でなければならない．

\sectionlink{与えられた FRW 解のまわりでの展開\label{app:tadpoles}}
本節では，与えられた FRW 背景（空間曲率に応じて $k=-1,0,+1$）のまわりで，時間微分同相変換が破れた最も一般的な理論を，
\be
S = \int d^4x \sqrt{-g} \left[\frac12 M_{\rm Pl}^2 R + M_{\rm Pl}^2 \Big(\dot H-\frac{k}{a^2}\Big) \cdot g^{00} - M_{\rm Pl}^2 \Big(3 H^2 +\dot H+2\frac{k}{a^2}\Big) + \ldots \right]
\label{eq:general}
\ee
と書けることを証明したい．ここで，省略記号は，空間微分同相変換のもとで不変であり，与えられた FRW 背景のまわりの揺らぎについて二次（またはそれ以上）の項を表す\footnote{$R$ の前の係数は，適切な場の再定義 $g_{\mu\nu} \to g_{\mu\nu} \cdot
  f(t)$ によって，常に時間に依存しないものにできる．これは通常の定式化では，Einsteinフレームへ移ることに対応する．}．

そのままでは，この主張は自明である．明示した項が望む FRW 発展を生むことはわかっているので，そこから離れないようにするには，追加の演算子はこの解のまわりで二次から始まらなければならない．ここで主張したいのは，追加される不変な項の{\em 一つ一つ}が，異なる演算子の間で一次の寄与を相殺させることなく，二次（またはそれ以上）となるということである．これらの項は，$g^{00}+1$，$\delta K_{\mu\nu}= K_{\mu\nu} -K_{\mu\nu}^{(0)}$，$\delta R_{\mu\nu\rho\sigma} \equiv R_{\mu\nu\rho\sigma} - R_{\mu\nu\rho\sigma}^{(0)}$ などの一次の演算子からなる，多項式（二次以上）として書かれる．与えられた FRW 解で評価した値を陽に差し引いているため，これらの演算子は摂動の一次から始まることに注意しよう．FRW 計量の対称性により，背景上で評価したあらゆるテンソル（$K_{\mu\nu}^{(0)}$，$R_{\mu\nu\rho\sigma}^{(0)}$，$(\nabla_\alpha R_{\mu\nu\rho\sigma})^{(0)}$ など）は，$g_{\mu\nu}$，$n_\mu$，および時間の関数だけを用いて書ける．たとえば，
\begin{eqnarray}
K^{(0)}_{\mu\nu} & =&  a^2 H h_{\mu\nu} \\
R_{\mu\nu\rho\sigma}^{(0)} & = & 2 (H+k) h_{\mu [\rho} h_{\sigma ] \nu} + (\dot H + H^2) a^2 h_{\mu\sigma} \delta_\nu^0 \delta_\rho^0 + {\rm perm.} 
\end{eqnarray}
となる．ここで，$k$ は空間的切片の曲率に依存する定数であり，置換は最後の項にのみ作用する．このように，FRW 背景上で評価したすべての演算子は，それ自体が共変な演算子である．したがって，$\delta K_{\mu\nu}$ や $\delta R_{\mu\nu\rho\sigma}$ のような演算子は，適切に定義された共変な演算子であり，与えられた FRW 背景上ではゼロとなり，摂動の一次から始まる．
背景上で評価したテンソルをこのように書き換えられるのは，FRW 背景の高い対称性のおかげに限られることを強調しておく．対称性の低い解，{\em たとえば}非一様な背景のまわりで展開する場合には，これは成り立たない．

式（\ref{eq:general}）では，明示した演算子だけが揺らぎの一次の項を含むため，$\sqrt{-g}\;g^{00}$ と $\sqrt{-g}$ の係数は，背景の FRW 解によって一意に決定される．

ラグランジアンを常に（\ref{eq:general}）の形にできることを見てみよう．二つのテンソル $T$ と $G$ の縮約からなる演算子をとると（より多くのテンソルへの一般化は容易である），
\be
TG = \delta T \delta G + T^{(0)} G + T G^{(0)} - T^{(0)} G^{(0)} \;. \label{eq:splitting}
\ee
と書ける．和の各項について考えよう．第 1 項は，望みどおり，摂動の二次から始まることが明示されている．すでに述べたように，FRW 背景の対称性により，非摂動のテンソル $T^{(0)}$ と $G^{(0)}$ は，$g_{\mu\nu}$，$n_\mu$，$t$ の関数として書ける．したがって，最後の項 $T^{(0)} G^{(0)}$ は，時間に依存する係数を持つ $g^{00}$ の多項式にすぎない．これは，$\sqrt{-g}\;g^{00}$ と $\sqrt{-g}$ の項に加えて，摂動の二次から始まることが明示された演算子を含む．
残るのは，$T^{(0)} G$ という形のテンソルである．これらの項も，式~（\ref{eq:general}）の一次の演算子と，揺らぎの二次から始まる演算子との和として書けることを証明したい．構成上，$G$ は $K_{\mu\nu}$ または $R_{\mu\nu\rho\sigma}$ のいずれかについて一次であり，それらに共変微分が作用したものである．共変微分は，繰り返し部分積分を行い，$T^{(0)}$ と演算子の時間に依存する係数に作用させることで処理できる．その際，外的曲率の項が生成されることがある．その場合は，共変微分がなくなるまで式~（\ref{eq:splitting}）を繰り返し適用すればよい\footnote{$g^{00}$ のべきが，$T^{(0)}$ から生じることもある．これらは $g^{00} = -1 +\delta g^{00}$ と書くことで処理でき，式~（\ref{eq:general}）の $g^{00}$ 演算子への追加の寄与と，摂動について陽に二次以上である項とが生成される．}．この結果，共変微分を含まない，可能なスカラーの一次の項だけが残る．すなわち，$K^{\mu}  {}_{\mu}$ と $R^{00}$ である．どちらも，より便利な形に書き換えられる．外的曲率の項は部分積分でき，
\begin{eqnarray}
&&\int d^4x \sqrt{-g}\; f(t) K^{\mu} {}_{\mu} =\int d^4x \sqrt{-g}\; f\; \nabla_{\mu} n^{\mu}\nonumber \\ &&\qquad=-\int d^4 x \sqrt{-g} \;n^\mu \partial_\mu f =\int d^4x \sqrt{-g} \sqrt{-g^{00}} \dot f \,. 
\end{eqnarray}
となる．一方，$R^{00}$ は，次の関係式 \cite{Wald:1984rg} を用いて処理できる．
\begin{eqnarray}
(-g^{00})^{-1} R^{00}=R_{\mu\nu}n^\mu n^\nu = K^2 -K_{\mu\nu}K^{\mu\nu}-\nabla_\mu(n^\mu\nabla_\nu n^\nu)+\nabla_\nu(n^\mu \nabla_\mu n^\nu)\, .
\end{eqnarray}
最後の二つの項は，再び部分積分できる．
\begin{eqnarray}
\int d^4x \sqrt{-g} f(t) \nabla_\mu(n^\mu\nabla_\nu n^\nu)=-\int d^4x \sqrt{-g}\; \partial_\mu f\; n^\mu K^\nu {}_\nu\, ,
\end{eqnarray}
\begin{eqnarray}
\int d^4 x \sqrt{-g} f(t) \nabla_\nu(n^\mu \nabla_\mu n^\nu)=-\int d^4 x \sqrt{-g}\; \partial_\nu f\; n^\mu\nabla_\mu n^\nu =0
\end{eqnarray}
最後の変形では，$\partial_\nu f \propto n_\nu$ を用いた．以上から，$K^\mu_\mu$ と $R^{00}$ は，式~（\ref{eq:general}）の一次の演算子と，揺らぎの二次から始まる不変な項とを用いて書けることが示された．

以上より，与えられた FRW 背景のまわりで時間微分同相変換が破れた理論の最も一般的なラグランジアンは，次の形に書けることを示した．
\begin{eqnarray}\label{eq:general2}
S &=& \int d^4x \sqrt{-g} \; \Big[\frac12 M_{\rm Pl}^2 R + M_{\rm Pl}^2 \Big(\dot H-\frac{k}{a^2}\Big) \cdot g^{00} - M_{\rm Pl}^2 \Big(3 H^2 +\dot H+2\frac{k}{a^2}\Big) +  \\ \nonumber &&  F^{(2)}(g^{00}+1,\delta K_{\mu\nu}, \delta R_{\mu\nu\rho\sigma};\nabla_\mu;t)\Big]
\end{eqnarray}
ここで，$F^{(2)}$ は，引数 $g^{00}+1$，$\delta K_{\mu\nu}$，$\delta R_{\mu\nu\rho\sigma}$ について二次から始まる．

